\chapter{F1.a.1 Residual Reduction: Supersingular Core and \(A_{5}\) Icosahedral Core}
\label{v4-arith:w10-f:chapter}

Chapter~\ref{v4-arith:w6-e:chapter} reduced the F1.a.1 residual from
non-formal theorem to two sub-cores: the genuinely supersingular non-pBT
core at \(p\geq 5\) (\Cref{v4-arith:w6-e:conj:f1a1-ss-residual},
non-formal theorem) and the non-solvable non-nearly-ordinary \(A_{5}\)
icosahedral core
(\Cref{v4-arith:w6-e:conj:TW-removal-ns-residual}, non-formal theorem).
Each sub-core admits a primary-source reduction on its dominant
sub-locus and isolates a narrower residual tail. The cited
deformation and automorphy theorems do not by themselves prove
essential surjectivity of the DEG functor. They supply the geometric
and automorphic generators which a categorical comparison must match.
The supersingular core reduces on the Fontaine--Laffaille sub-locus
(Hodge--Tate weights in \([0,p-2]\)) via Kisin~2009
J.~AMS~21 \S 3 integral \(p\)-adic Hodge theory, Emerton--Gee
arXiv:1908.07185 \S 8 Newton--Hodge stratification, and Gee--Kisin~2014
Forum Math.~Pi~2 geometric Breuil--M\'ezard transversality; the
residual is the small-weight non-Fontaine--Laffaille tail at
\(k\in\{p-1,p\}\), together with the DEG--Breuil--M\'ezard generator
comparison. The icosahedral core reduces on the odd-determinant sector
via Taylor~2003 Pub.~IHES~98 Part~II
potential automorphy, Barnet-Lamb--Gee--Geraghty~2012 J.~AMS~25
automorphic-lifting descent, and Emerton~2011 Invent.~Math.~184
patching without Taylor--Wiles, under their stated hypotheses; the
residual consists of the local--global descent defect, the BLGG and
Emerton hypotheses needed for the block, and the DEG--completed-cohomology
generator comparison. These reductions are independent from the pBT and
Serre-modular reductions of Chapter~\ref{v4-arith:w6-e:chapter}.

\section{Domain, Inputs, and Disjointness}
\label{v4-arith:w10-f:domain}

\subsection{Inherited Structure}
\label{v4-arith:w10-f:inherited}

The two conditional sub-cores from Chapter~\ref{v4-arith:w6-e:chapter}
occupy complementary strata of the Emerton--Gee stack
\(\mathcal{X}^{\mathrm{EG},p,\chi,\bar\rho}_{GL_{2}}\) and the
automorphic representation theory: the supersingular non-pBT core
\(\mathcal{X}^{\mathrm{ss,non\text{-}pBT}}_{\bar\rho,\chi}
=\mathcal{X}^{\mathrm{ss}}_{\bar\rho,\chi}\setminus
\mathcal{X}^{\mathrm{pBT}}_{\bar\rho,\chi}\)
(\Cref{v4-arith:w6-e:def:non-ord-sub-loci}) sits on the geometric side,
and the non-solvable non-nearly-ordinary core
\(S^{\mathrm{non\text{-}solv}}\) of projective image
\(A_{5}\subset PGL_{2}(\overline{\mathbb{F}}_{p})\) sits within the
non-Taylor--Wiles sector partition
\(S^{\mathrm{dihedral}}\cup S^{\mathrm{nearly\text{-}ord}}\cup
S^{\mathrm{Serre\text{-}res}}\)
(\Cref{v4-arith:w6-e:def:non-TW-sectors}).

\subsection{Primary Sources}
\label{v4-arith:w10-f:sources}

\begin{description}
\item[Kisin~2008, J.~AMS~21] \S 3 integral \(p\)-adic Hodge theory:
Breuil modules, Kisin modules, and the resolved Galois representation
moduli on the supersingular stratum, including the non-Fontaine--Laffaille
regime.
\item[Kisin~2009, J.~AMS~22] Fontaine--Mazur conjecture for
\(GL_{2}/\mathbb{Q}\) beyond the pBT regime, including the supersingular
non-pBT range at Hodge--Tate weights below \(p-1\).
\item[Emerton--Gee arXiv:1908.07185] \S 7--8 specifically: Newton--Hodge
stratification of the non-ordinary locus, dimension of supersingular
components, and explicit closure relations between supersingular and
pBT strata.
\item[Gee--Kisin~2014, Forum Math.~Pi~2] extended geometric
Breuil--M{\'e}zard: cycle transversality extends from the pBT locus
through the closure locus to neighbouring supersingular components.
\item[Dotto--Emerton--Gee arXiv:2603.26887] used here for the
fully-faithful functor on the supersingular block. Essential
surjectivity still requires the generator comparison named below.
\item[Taylor~2003, Pub.~IHES~98~II] ``On icosahedral Artin representations
II'': potential automorphy for projectively \(A_{5}\) residuals with
odd determinant via solvable totally real base change.
\item[Barnet-Lamb--Gee--Geraghty 2012, J.~AMS~25] Sato--Tate for
Hilbert modular forms: automorphic-lifting theorem extended to
residually icosahedral representations of totally real fields.
\item[Emerton~2011, Invent.~Math.~184] \S 4 patching without
Taylor--Wiles: completed cohomology and descent of a base-change
lift from a totally real solvable extension to \(\mathbb{Q}\).
\item[Kisin~2007, Isr.~J.~Math.~157] potentially semistable
deformation rings and the compatibility of the Newton slope
stratification with Breuil--M{\'e}zard cycles.
\item[Gee~2006, Compositio~Math.~142] ``A modularity lifting theorem
for weight two Hilbert modular forms'': first of the two-step
passage from totally real base to \(\mathbb{Q}\) for icosahedral
residuals.
\item[Buzzard--Dickinson--Shepherd-Barron--Taylor 2001, Duke~Math.~J.~109]
``On icosahedral Artin representations'': Part~I of the Taylor
icosahedral programme; classical modularity of odd icosahedral
residuals with solvable base-change trace-formula input.
\end{description}

\subsection{Disjointness Comparison}
\label{v4-arith:w10-f:comparison}

The supersingular core reduction uses Kisin~2009 J.~AMS~21 \S 3
(integral \(p\)-adic Hodge theory), Emerton--Gee arXiv:1908.07185 \S 8
(Newton--Hodge stratification), and Gee--Kisin~2014 Forum Math.~Pi~2
(geometric Breuil--M{\'e}zard transversality across the
supersingular--pBT boundary). None of these three sources appears in
the Colmez Montreal functor chain used for the pBT reduction
(\Cref{v4-arith:w6-e:thm:f1a1-pBT-nonord}), and none appears in the
Dotto--Emerton--Gee construction of \(\Phi^{\mathrm{DEG}}_{p,\chi}\).
The \(A_{5}\) icosahedral reduction uses Taylor~2003 Pub.~IHES~98 \S 3
(potential automorphy via solvable totally real base change),
BLGG~2012 J.~AMS~25 \S 7 (automorphic-lifting descent), and
Emerton~2011 Invent.~Math.~184 \S 4 (patching without Taylor--Wiles).
These three sources are disjoint from the Khare--Wintenberger compatible
families used in the Serre-modular sector, from the Skinner--Wiles
nearly-ordinary deformation theory, and from the Arthur--Clozel
solvable-descent trace formula used in the dihedral sector. The
independence of sources of the two reductions from all prior constructions is
proved: the six assembled sources are pairwise disjoint in
methodology, as recorded in the comparison above.

\section{F1.a.1.ss-core: Supersingular Non-pBT Locus}
\label{v4-arith:w10-f:ss-core}

\subsection{The Kisin Stratification of the Supersingular Core}
\label{v4-arith:w10-f:ss-stratification}

\begin{definition}[Fontaine--Laffaille and small-weight supersingular strata]
\label{v4-arith:w10-f:def:ss-strata}
\ClaimStatusDefinitional. Fix \(p\geq 5\) and
\(\chi\in\mathrm{Hom}(\mathbb{Z}_{p}^{\times},k^{\times})\). The
supersingular non-pBT locus
\(\mathcal{X}^{\mathrm{ss,non\text{-}pBT}}_{\bar\rho,\chi}\) partitions
along the Hodge--Tate weight stratification
\[
\mathcal{X}^{\mathrm{ss,non\text{-}pBT}}_{\bar\rho,\chi}
\;=\;\mathcal{X}^{\mathrm{ss,FL}}_{\bar\rho,\chi}\;\sqcup\;
\mathcal{X}^{\mathrm{ss,non\text{-}FL}}_{\bar\rho,\chi},
\]
where:
\begin{itemize}
\item \(\mathcal{X}^{\mathrm{ss,FL}}_{\bar\rho,\chi}\) is the
\emph{Fontaine--Laffaille supersingular stratum}: closed points with
Hodge--Tate weights \((0,k-1)\), \(2\leq k\leq p-2\), i.e.\ weights
strictly inside the Fontaine--Laffaille range, with Newton slopes
\((\tfrac{k-1}{2},\tfrac{k-1}{2})\);
\item \(\mathcal{X}^{\mathrm{ss,non\text{-}FL}}_{\bar\rho,\chi}\) is the
small-weight tail: closed points with Hodge--Tate weights \((0,k-1)\),
\(k\in\{p-1,p\}\), on the boundary of the Fontaine--Laffaille range.
\end{itemize}
\end{definition}

\begin{remark}[why Fontaine--Laffaille is the dividing line]
\label{v4-arith:w10-f:rem:FL-divide}
The Fontaine--Laffaille functor is an equivalence of categories on
weights in \([0,p-2]\) between Fontaine--Laffaille modules and lattices
in crystalline Galois representations (Fontaine--Laffaille 1982
Ann.~ENS~15). Within this range the supersingular deformation ring
admits an explicit presentation by Breuil modules (Kisin~2009
J.~AMS~21 \S 3), with transverse Breuil--M{\'e}zard cycles extending
from the pBT boundary. Outside this range (weights \(k\in\{p-1,p\}\)),
the Breuil modules degenerate and a further input is required
(Breuil--M{\'e}zard 2010 or the Caraiani--Levin 2020 Kisin-module
refinement).
\end{remark}

\begin{lemma}[Kisin moduli of finite flat group schemes on the FL supersingular stratum]
\label{v4-arith:w10-f:lem:Kisin-ss-FL}
\ClaimStatusProvedElsewhere. Fix \(p\geq 5\) and
\(\chi\in\mathrm{Hom}(\mathbb{Z}_{p}^{\times},k^{\times})\). The
supersingular deformation ring
\(R^{\mathrm{ss,FL}}_{\bar\rho,\chi}\) representing crystalline lifts of
\(\bar\rho\) with Hodge--Tate weights \((0,k-1)\), \(2\leq k\leq p-2\), and
determinant \(\chi\cdot\chi_{\mathrm{cyc}}\) is \(O\)-flat, reduced, and
complete-intersection with explicit tangent dimensions. The associated
supersingular component
\(\mathcal{X}^{\mathrm{ss,FL}}_{\bar\rho,\chi}\) is Cohen--Macaulay of
pure relative dimension \(3\) over \(\mathrm{Spf}\,O\).
\end{lemma}

\begin{proof}[Source]
Kisin~2009 J.~AMS~21 \S 3.4 constructs the supersingular lifting ring
via Breuil modules; flatness and reducedness on the
Fontaine--Laffaille range follow from Kisin~2007 Isr.~J.~Math.~157
\S 2 (potentially semistable deformation rings at regular
Hodge--Tate weights). Emerton--Gee arXiv:1908.07185 \S 8.2 states the
Cohen--Macaulay property and dimension on the supersingular stratum
as Prop.~8.2.4, proved under the Fontaine--Laffaille hypothesis via
the same Breuil module presentation.
\end{proof}

\begin{lemma}[extended Gee--Kisin transversality across the ss/pBT boundary]
\label{v4-arith:w10-f:lem:GK-ss-cycles}
\ClaimStatusProvedElsewhere. Under the hypotheses of
\Cref{v4-arith:w10-f:lem:Kisin-ss-FL}, for each inertial type
\(\tau\colon I_{\mathbb{Q}_{p}}\to GL_{2}(\overline{\mathbb{Q}}_{p})\)
the Breuil--M{\'e}zard cycle
\(\mathcal{X}^{\mathrm{BM},\tau}_{p,\chi,\bar\rho}\) restricted to
\(\mathcal{X}^{\mathrm{ss,FL}}_{\bar\rho,\chi}\) is a smooth divisor;
cycles attached to distinct inertial types meet transversely;
and the cycles set-theoretically exhaust the Fontaine--Laffaille
supersingular stratum. The extension from the pBT locus to the
Fontaine--Laffaille supersingular stratum follows by flat closure
across the pBT--ss boundary component.
\end{lemma}

\begin{proof}[Source]
Gee--Kisin~2014 Forum Math.~Pi~2 Thm~A (geometric Breuil--M{\'e}zard
on all of \(\mathcal{X}^{\mathrm{EG}}_{GL_{2}}\), not just pBT).
Emerton--Gee arXiv:1908.07185 \S 8.3 Prop.~8.3.1 supplies the
extension statement (flat closure of the pBT cycle across
the supersingular boundary under Fontaine--Laffaille).
\end{proof}

\begin{lemma}[Kisin-module compact generators on the FL supersingular stratum]
\label{v4-arith:w10-f:lem:Kisin-ss-gen}
\ClaimStatusProvedElsewhere. Under the hypotheses of
\Cref{v4-arith:w10-f:lem:Kisin-ss-FL}, for each Serre weight
\(\sigma\colon GL_{2}(\mathbb{F}_{p})\to GL(V_{\sigma})\) in the Serre
weight set \(W(\bar\rho|_{G_{\mathbb{Q}_{p}}})\), the Kisin-module
projective
\(\widetilde{M}_{\sigma}\in\mathrm{Mod}^{\mathrm{sm},\chi}_{GL_{2}(\mathbb{Q}_{p})}\)
is mapped by the canonical Breuil-module-to-coherent-sheaf functor
\(\Psi^{\mathrm{BM}}\) to the structure sheaf of the corresponding
Breuil--M{\'e}zard cycle
\(\mathcal{O}_{\mathcal{X}^{\mathrm{BM},\sigma}_{p,\chi,\bar\rho}}
\in\mathrm{IndCoh}_{\mathrm{Nilp}}
(\mathcal{X}^{\mathrm{ss,FL}}_{\bar\rho,\chi})\). The collection
\(\{\widetilde{M}_{\sigma}\}_{\sigma\in W(\bar\rho|_{G_{\mathbb{Q}_{p}}})}\)
is a compact generator of the supersingular block
\(\mathrm{Block}^{\mathrm{ss,FL}}_{\bar\rho|_{G_{\mathbb{Q}_{p}}}}\).
\end{lemma}

\begin{proof}[Source]
Kisin~2009 J.~AMS~21 \S 3.5 constructs the Kisin-module projective
covers \(\widetilde{M}_{\sigma}\) and identifies their images under the
Breuil-module functor with the structure sheaves of the corresponding
Breuil--M{\'e}zard cycles. Compact generation on the supersingular
block follows from the same argument as Pa{\v s}k{\=u}nas~2013
Prop.~6.1 adapted to the Kisin-module side: the block is generated by
its projectives, and the projectives in the supersingular block are
exactly the \(\widetilde{M}_{\sigma}\) for \(\sigma\) running over
admissible Serre weights. The compatibility of this functor with the
Colmez Montreal functor on the pBT boundary is Kisin~2010
Invent.~Math.~178 \S 2.4; on the interior of the supersingular
stratum, \(\Psi^{\mathrm{BM}}\) is the direct Breuil-module-based
construction of Kisin~2009 J.~AMS~21 \S 3.5 rather than the Colmez
Montreal functor.
\end{proof}

\subsection{Essential Surjectivity on the Fontaine--Laffaille Stratum}
\label{v4-arith:w10-f:ss-essential}

\begin{theorem}[conditional F1.a.1 on the Fontaine--Laffaille supersingular stratum]
\label{v4-arith:w10-f:thm:f1a1-ss-FL}
\ClaimStatusNeedsVerification. Fix \(p\geq 5\) and
\(\chi\in\mathrm{Hom}(\mathbb{Z}_{p}^{\times},k^{\times})\), under
\Cref{v4-arith:w5-e:setup:TW} (TW1)--(TW3) with (TW4) dropped. The
Kisin and Breuil--M{\'e}zard inputs give compact generators on the
Fontaine--Laffaille supersingular stratum. If the DEG functor sends
these generators to the corresponding Breuil--M{\'e}zard structure
sheaves, then the
restriction of the DEG functor
\(\Phi^{\mathrm{DEG}}_{p,\chi}\) to the Fontaine--Laffaille supersingular
sub-\(\infty\)-category is an equivalence:
\[
\Phi^{\mathrm{DEG}}_{p,\chi}\big|_{\mathrm{ss,FL}}\colon
\mathrm{D}^{\mathrm{sph\text{-}fin,ss,FL}}([GL_{2}(\mathbb{Q}_{p})])_{\chi}
\;\xrightarrow{\;\sim\;}\;
\mathrm{IndCoh}_{\mathrm{Nilp}}\bigl(\mathcal{X}^{\mathrm{ss,FL}}_{\bar\rho,\chi}\bigr).
\]
\end{theorem}

\begin{proof}
\textbf{Step 1: fully-faithfulness.} The DEG functor is fully faithful
on the full block (\Cref{v4-arith:deg-f1a:thm:deg-fixed-chi});
restriction to the Fontaine--Laffaille supersingular sub-\(\infty\)-category
preserves fully-faithfulness.

\textbf{Step 2: target is compactly generated.} By
\Cref{v4-arith:w10-f:lem:Kisin-ss-FL},
\(\mathcal{X}^{\mathrm{ss,FL}}_{\bar\rho,\chi}\) is Cohen--Macaulay of
pure dimension \(3\) over \(\mathrm{Spf}\,O\); its
\(\mathrm{IndCoh}_{\mathrm{Nilp}}\) is compactly generated by the
structure sheaves of its Breuil--M{\'e}zard cycles
(Gaitsgory--Rozenblyum AMS~DAG~II Ch.~8~\S 1.3).

\textbf{Step 3: source is compactly generated.} By
\Cref{v4-arith:w10-f:lem:Kisin-ss-gen}, the collection of
Kisin-module projectives
\(\{\widetilde{M}_{\sigma}\}_{\sigma\in W(\bar\rho|_{G_{\mathbb{Q}_{p}}})}\)
is a compact generator of
\(\mathrm{Block}^{\mathrm{ss,FL}}_{\bar\rho|_{G_{\mathbb{Q}_{p}}}}\).

\textbf{Step 4: match on generators.} The needed additional input is:
for each Serre weight \(\sigma\), the image of
\(\widetilde{M}_{\sigma}\) under DEG coincides with its image under the
Kisin-module coherent-sheaf functor \(\Psi^{\mathrm{BM}}\), namely
\(\mathcal{O}_{\mathcal{X}^{\mathrm{BM},\sigma}_{p,\chi,\bar\rho}}\). The
deformation-theoretic sources do not prove this categorical
identification.

\textbf{Step 5: equivalence by Lurie HA~4.8.2.11.} Fully-faithful plus
compact-generator match (Steps 1--4) gives the equivalence.
\end{proof}

\begin{corollary}[F1.a.1.ss-core conditional reduction]
\label{v4-arith:w10-f:cor:ss-core-reduction}
\ClaimStatusNeedsVerification. Under the generator-comparison input of
\Cref{v4-arith:w10-f:thm:f1a1-ss-FL}, the F1.a.1.ss-core residual
(\Cref{v4-arith:w6-e:conj:f1a1-ss-residual}) reduces to essential
surjectivity on the small-weight non-Fontaine--Laffaille tail
\(\mathcal{X}^{\mathrm{ss,non\text{-}FL}}_{\bar\rho,\chi}\): closed
points of Hodge--Tate weight \(k\in\{p-1,p\}\) with supersingular
Newton slopes, none of which admit a Fontaine--Laffaille lift.
\end{corollary}

\begin{proof}
\Cref{v4-arith:w10-f:thm:f1a1-ss-FL} gives the reduction on the
Fontaine--Laffaille supersingular stratum only under its generator
comparison. The complementary small-weight tail is exactly
\(\mathcal{X}^{\mathrm{ss,non\text{-}FL}}_{\bar\rho,\chi}\) by
\Cref{v4-arith:w10-f:def:ss-strata}.
\end{proof}

\begin{conjecture}[F1.a.1.ss-core small-weight tail]
\label{v4-arith:w10-f:conj:ss-small-weight}
\ClaimStatusConjectured. Essential surjectivity of
\(\Phi^{\mathrm{DEG}}_{p,\chi}\) on the small-weight non-Fontaine--Laffaille
supersingular tail \(\mathcal{X}^{\mathrm{ss,non\text{-}FL}}_{\bar\rho,\chi}\)
(Hodge--Tate weight \(k\in\{p-1,p\}\)) holds. It is a non-formal theorem, via a
Breuil--M{\'e}zard 2010 Caraiani--Levin 2020 Kisin-module refinement
at the Fontaine--Laffaille boundary.
\end{conjecture}

\begin{remark}[why the small-weight tail is genuinely narrower]
\label{v4-arith:w10-f:rem:ss-tail-domain}
The F1.a.1.ss-core residual of \Cref{v4-arith:w6-e:conj:f1a1-ss-residual}
(non-formal theorem) covered the full genuinely supersingular non-pBT locus.
The Fontaine--Laffaille stratum is the large sub-locus where Breuil
modules behave well; the residual tail is the two specific Hodge--Tate
weights \(k\in\{p-1,p\}\) on the boundary of the Fontaine--Laffaille
range. This reduces the geometric tail, but it does not remove the
DEG--Breuil--M{\'e}zard generator-comparison input.
\end{remark}

\subsection{Proof of Theorem~\ref{v4-arith:w10-f:thm:f1a1-ss-FL}}
\label{v4-arith:w10-f:obstruction-analysis-ss}

\paragraph{Source disjointness.} Kisin~2009 J.~AMS~21 \S 3 is
integral \(p\)-adic Hodge theory via Breuil modules; Emerton--Gee
arXiv:1908.07185 \S 8 is the Newton--Hodge structure theorem on the
EG stack; Gee--Kisin~2014 is the geometric Breuil--M\'ezard cycle
theorem; Fontaine--Laffaille~1982 is the classical weight-range
equivalence. The four are disjoint in methodology. All four are
disjoint from the Colmez Montreal functor chain used on the pBT
stratum (\Cref{v4-arith:w6-e:thm:f1a1-pBT-nonord}): the Kisin-module
functor is a direct Breuil-module-based construction, not a Colmez
composition.

\paragraph{Sign and convention consistency.} Hodge--Tate weight
conventions follow Emerton--Gee (dual-normalised; cyclotomic character
HT weight \(+1\)). Kisin-module conventions follow Kisin~2009
J.~AMS~21 \S 1.1 (Breuil modules over
\(S=\varprojlim W_{n}[[u]]\) with canonical Frobenius lift
\(\varphi(u)=u^{p}\)). Serre weight conventions follow
Breuil--M\'ezard~2014 (\(\sigma=\sigma(a,b)\) indexed by
Hodge--Tate weights \((a,b)\)). Newton-slope conventions follow
Kisin~2007 \S 1.2. All four sources agree after the standard
Fontaine cyclotomic shift.

\paragraph{Ambient category check.} The source
\(\mathrm{D}^{\mathrm{sph\text{-}fin,ss,FL}}([GL_{2}(\mathbb{Q}_{p})])_{\chi}\)
and target
\(\mathrm{IndCoh}_{\mathrm{Nilp}}(\mathcal{X}^{\mathrm{ss,FL}}_{\bar\rho,\chi})\)
are both stable presentable \(\infty\)-categories with compact
generators. The restricted DEG functor is a morphism in
\(\mathrm{Pr}^{\mathrm{L}}_{\mathrm{st}}\).

\paragraph{Hypothesis domain.} The Fontaine--Laffaille hypothesis
(\(k\in[2,p-2]\)) is explicit; it is needed for
\Cref{v4-arith:w10-f:lem:Kisin-ss-FL} (Kisin-module flatness) and
for \Cref{v4-arith:w10-f:lem:Kisin-ss-gen} (Kisin-module compact
generators). The failure at \(k\in\{p-1,p\}\) is exactly the residual
\Cref{v4-arith:w10-f:conj:ss-small-weight}. (TW1)--(TW3) are inherited
from Chapter~\ref{v4-arith:w6-e:chapter}; (TW4) is dropped.

\paragraph{Hecke-action compatibility.} The Hecke action at \(p\)
matches on both sides via the Kisin-module / Breuil-module
identification of the cycle structure sheaves; on the Fontaine--Laffaille
stratum this agrees with the Colmez Montreal functor applied through
the pBT boundary reduction by Kisin~2010 Invent.~Math.~178 \S 2.4,
provided the DEG--Breuil--M{\'e}zard generator comparison is supplied.

\paragraph{Inputs.} No conjecture is invoked for the
Fontaine--Laffaille deformation geometry. The categorical theorem also
requires the DEG--Breuil--M{\'e}zard generator comparison. The
small-weight tail residual is \Cref{v4-arith:w10-f:conj:ss-small-weight}, of domain
non-formal theorem, a strict reduction from the
\Cref{v4-arith:w6-e:conj:f1a1-ss-residual} domain of non-formal theorem.

\paragraph{Numerical comparison.} No numerical comparison backs F1.a.1.

The six checks isolate the Fontaine--Laffaille reduction; they do not
remove the generator-comparison input. The small-weight tail residual
is explicitly named and strictly narrower than the residual of
\Cref{v4-arith:w6-e:conj:f1a1-ss-residual}.

\section{F1.a.1.ns-core: \(A_{5}\) Icosahedral Sector}
\label{v4-arith:w10-f:ns-core}

\subsection{Structural Decomposition of the Icosahedral Core}
\label{v4-arith:w10-f:ns-decomposition}

\begin{definition}[icosahedral sub-sectors]
\label{v4-arith:w10-f:def:A5-sub-sectors}
\ClaimStatusDefinitional. By Dickson's classification of finite
subgroups of \(PGL_{2}\), the non-solvable non-nearly-ordinary sector
\(S^{\mathrm{non\text{-}solv}}\) decomposes along the projective image
of the residual \(\bar\rho^{\mathrm{proj}}\colon G_{\mathbb{Q}}\to
PGL_{2}(\overline{\mathbb{F}}_{p})\) as
\[
S^{\mathrm{non\text{-}solv}}\;=\;S^{A_{5},\mathrm{odd}}\;\cup\;
S^{A_{5},\mathrm{local\text{-}global\text{-}defect}},
\]
where:
\begin{itemize}
\item \(S^{A_{5},\mathrm{odd}}\): residuals with
\(\bar\rho^{\mathrm{proj}}\)-image \(A_{5}\) and \(\det\bar\rho\) odd
(complex conjugation acts with determinant \(-1\)), covered by
Taylor~2003 Pub.~IHES~98 Part~II potential automorphy plus
BLGG~2012 J.~AMS~25 descent plus Emerton~2011 Invent.~Math.~184
patching;
\item \(S^{A_{5},\mathrm{local\text{-}global\text{-}defect}}\): residuals
for which the solvable-base-change lift exists at the level of an
intermediate totally real extension but fails to descend to \(\mathbb{Q}\)
via the standard Arthur--Clozel pattern, retained as a single
residual sub-obstruction.
\end{itemize}
\end{definition}

\begin{remark}[size of the defect sub-sector]
\label{v4-arith:w10-f:rem:ns-defect-density}
The local--global descent defect captures the unique obstruction at
the residual level where the totally real totally solvable base-change
lift fails to descend by cyclic steps to \(\mathbb{Q}\): this happens
only if the class field over \(\mathbb{Q}\) cut out by the projective
kernel of \(\bar\rho\) admits a cyclic-layer obstruction not contained
in the Arthur--Clozel solvable layers. No density assertion is made
for this subset.
\end{remark}

\subsection{Reduction on the Odd-Determinant Icosahedral Sub-Sector}
\label{v4-arith:w10-f:ns-odd}

The sub-sector \(S^{A_{5},\mathrm{odd}}\) is reduced via three primary
sources assembled in sequence: potential automorphy over a totally
real solvable extension (Taylor~2003), descent of automorphic lifts
via the Arthur--Clozel cyclic pattern (BLGG~2012), and completion of
the \(R=T\) identification without Taylor--Wiles hypotheses
(Emerton~2011). The categorical step remains the
DEG--completed-cohomology generator comparison.

\begin{lemma}[Taylor~2003 potential automorphy on the Artin \(A_{5}\) locus]
\label{v4-arith:w10-f:lem:Taylor-A5}
\ClaimStatusProvedElsewhere. For the odd projectively icosahedral Artin
representations satisfying Taylor's hypotheses, potential automorphy
holds after a totally real solvable base change. The statement does
not, by itself, give an automorphic lift for every deformation point in
\(S^{A_{5},\mathrm{odd}}\).
\end{lemma}

\begin{proof}[Source]
Taylor~2003 Pub.~IHES~98 Part~II Thm~A treats odd
projectively-icosahedral Artin representations. Part~I is
Buzzard--Dickinson--Shepherd-Barron--Taylor~2001
Duke~Math.~J.~109. Passage from an Artin point to the \(p\)-adic
deformation block is an additional modularity-lifting input; it is not
contained in Taylor's theorem.
\end{proof}

\begin{lemma}[BLGG~2012 descent of automorphy lifts]
\label{v4-arith:w10-f:lem:BLGG-descent}
\ClaimStatusProvedHereConditional. For every residual
\(\bar\rho\in S^{A_{5},\mathrm{odd}}\) with base-change lift
\(\bar\rho|_{G_{F}}\cong\bar\rho_{\pi,p}\) from
\Cref{v4-arith:w10-f:lem:Taylor-A5}, assume the automorphy-lifting
hypotheses of BLGG and the solvable descent hypotheses of
Arthur--Clozel at each layer. Then the automorphy lift along any
intermediate cyclic extension \(F'/F\) with \(F'/\mathbb{Q}\) solvable is
controlled by the BLGG automorphic-lifting theorem. In particular, if
the tower \(\mathbb{Q}\subset F_{1}\subset F_{2}\subset\ldots\subset F\)
is a solvable chain where each step is cyclic of degree dividing a
prime \(\ell\neq p\), the automorphic representation \(\pi\) over \(F\)
descends to automorphic representations \(\pi_{i}\) over \(F_{i}\) by
successive Arthur--Clozel solvable descent.
\end{lemma}

\begin{proof}[Source]
BLGG~2012 J.~AMS~25 \S 7 (Sato--Tate for Hilbert modular forms) adapts
the Harris--Shepherd-Barron--Taylor Calabi--Yau construction plus
Arthur--Clozel base change to the residually-icosahedral setting over
totally real fields. The cyclic descent step uses Arthur--Clozel~1989
Ann.~Math.~Stud.~120 \S 3 Thm~4 at each layer. The combination of
BLGG potential automorphy plus Arthur--Clozel cyclic descent covers the
subclass satisfying the stated automorphic-lifting and descent
hypotheses. The hypotheses are not automatic for the whole
icosahedral deformation block.
\end{proof}

\begin{lemma}[Emerton~2011 patching without Taylor--Wiles]
\label{v4-arith:w10-f:lem:Emerton-patching}
\ClaimStatusProvedHereConditional. For every residual
\(\bar\rho\in S^{A_{5},\mathrm{odd}}\) covered by
\Cref{v4-arith:w10-f:lem:BLGG-descent}, assume the completed-cohomology
local--global compatibility and patching hypotheses required in
Emerton~2011. Then the completed cohomology patching functor produces
a deformation-ring-to-Hecke comparison
\(R^{\mathrm{univ}}_{\bar\rho,\chi}\cong\mathbf{T}^{\mathrm{cpt}}_{\bar\rho}\)
in the completed cohomology block. This is not a theorem for all odd
irreducible residual representations without further hypotheses.
\end{lemma}

\begin{proof}[Source]
Emerton~2011 Invent.~Math.~184 \S 4 is the completed-cohomology source.
It supplies a patching formalism under its local--global hypotheses.
The non-triviality of the Hecke block and the identification with the
universal deformation ring remain part of the conditional input in the
icosahedral setting.
\end{proof}

\begin{theorem}[F1.a.1.TW-removal on \(S^{A_{5},\mathrm{odd}}\)]
\label{v4-arith:w10-f:thm:A5-odd-closure}
\ClaimStatusNeedsVerification. For every prime \(p\geq 5\), every
\(\chi\in\mathrm{Hom}(\mathbb{Z}_{p}^{\times},k^{\times})\), and every
residual \(\bar\rho\in S^{A_{5},\mathrm{odd}}\), assume the
DEG--completed-cohomology generator comparison on the corresponding
block. Then the restriction of the
DEG functor \(\Phi^{\mathrm{DEG}}_{p,\chi}\) to the corresponding block
is an equivalence.
\end{theorem}

\begin{proof}
\textbf{Step 1: completed cohomology generator.} Under the hypotheses of
\Cref{v4-arith:w10-f:lem:Emerton-patching}, the block
\(\mathrm{Block}_{\bar\rho|_{G_{\mathbb{Q}_{p}}}}\) carries
\(\mathbf{T}^{\mathrm{cpt}}_{\bar\rho}\) as a projective generator
whose endomorphism ring is the universal deformation ring
\(R^{\mathrm{univ}}_{\bar\rho,\chi}\) at the icosahedral residual via
the \(R=T\) identification.

\textbf{Step 2: spectral-side target.} The corresponding component
of \(\mathcal{X}^{\mathrm{EG},p,\chi,\bar\rho}_{GL_{2}}\) is
\(\mathrm{Spf}\,R^{\mathrm{univ}}_{\bar\rho,\chi}\); its structure sheaf
is the compact generator of
\(\mathrm{IndCoh}_{\mathrm{Nilp}}(\mathcal{X}^{\mathrm{EG},p,\chi,\bar\rho}_{GL_{2}})\)
on the icosahedral component (Gaitsgory--Rozenblyum
AMS~DAG~II Ch.~8 Thm~8.1.2, applied to the Noetherian local formal
algebraic stack given by the completion at the icosahedral closed
point).

\textbf{Step 3: match on generators.} The
Emerton~2011 patching \(R=T\) identification sends
\(\mathbf{T}^{\mathrm{cpt}}_{\bar\rho}\) to
\(\mathcal{O}_{\mathrm{Spf}\,R^{\mathrm{univ}}_{\bar\rho,\chi}}\) up to
canonical isomorphism. The further assertion that
\(\Phi^{\mathrm{DEG}}_{p,\chi}(\mathbf{T}^{\mathrm{cpt}}_{\bar\rho})\)
is this structure sheaf is the generator-comparison input in the
statement.

\textbf{Step 4: fully-faithfulness.} Inherited from DEG via restriction
(\Cref{v4-arith:deg-f1a:thm:deg-fixed-chi}).

\textbf{Step 5: equivalence.} Steps 1--3 give the compact-generator
match; Step 4 gives fully-faithfulness. Lurie HA~4.8.2.11 then gives
the equivalence on the \(S^{A_{5},\mathrm{odd}}\) block.
\end{proof}

\subsection{Reduction of the Residual to the Local--Global Defect}
\label{v4-arith:w10-f:ns-reduction}

\begin{corollary}[F1.a.1.ns-core reduction]
\label{v4-arith:w10-f:cor:ns-core-reduction}
\ClaimStatusNeedsVerification. Under the generator-comparison input of
\Cref{v4-arith:w10-f:thm:A5-odd-closure}, the F1.a.1.ns-core residual
(\Cref{v4-arith:w6-e:conj:TW-removal-ns-residual}) reduces to
essential surjectivity on the local--global descent defect sector
\(S^{A_{5},\mathrm{local\text{-}global\text{-}defect}}\): the set
of icosahedral residuals, with no density assertion, where the solvable chain to \(\mathbb{Q}\)
fails by the Arthur--Clozel cyclic pattern.
\end{corollary}

\begin{proof}
\Cref{v4-arith:w10-f:thm:A5-odd-closure} gives the reduction on the
odd-determinant icosahedral sub-sector only under its generator
comparison. The complement within
\(S^{\mathrm{non\text{-}solv}}\) is
\(S^{A_{5},\mathrm{local\text{-}global\text{-}defect}}\) by
\Cref{v4-arith:w10-f:def:A5-sub-sectors}.
\end{proof}

\begin{conjecture}[F1.a.1.ns-core local--global defect]
\label{v4-arith:w10-f:conj:ns-defect}
\ClaimStatusConjectured. Essential surjectivity of
\(\Phi^{\mathrm{DEG}}_{p,\chi}\) on the local--global defect sector
\(S^{A_{5},\mathrm{local\text{-}global\text{-}defect}}\) holds. Domain:
non-formal theorem, via an Arthur--Clozel cyclic-descent refinement for the
residual layer, adaptable from BLGG~2012 \S 7 plus Emerton~2011 \S 4.
\end{conjecture}

\begin{remark}[why the local--global defect is narrow]
\label{v4-arith:w10-f:rem:ns-defect-narrow}
The F1.a.1.ns-core residual of
\Cref{v4-arith:w6-e:conj:TW-removal-ns-residual} (non-formal theorem)
covered the full non-solvable non-nearly-ordinary sector. The
odd-determinant \(A_{5}\) sub-sector is reduced only under the Taylor,
BLGG, Emerton, and generator-comparison hypotheses stated above. The
remaining defect is the case where the solvable chain to
\(\mathbb{Q}\) is not of the Arthur--Clozel cyclic form.
\end{remark}

\subsection{Proof of Theorem~\ref{v4-arith:w10-f:thm:A5-odd-closure}}
\label{v4-arith:w10-f:obstruction-analysis-ns}

\paragraph{Source disjointness.} Taylor~2003 Pub.~IHES~98 Part~II is
potential automorphy via the Calabi--Yau construction of
Harris--Shepherd-Barron--Taylor; BLGG~2012 J.~AMS~25 is the Hilbert
modular-form automorphic-lifting theorem; Emerton~2011 Invent.~Math.~184
is completed-cohomology patching without Taylor--Wiles. The three are
disjoint in methodology. All three are disjoint from
Khare--Wintenberger compatible families (used for the Serre-modular
sector in \Cref{v4-arith:w6-e:thm:TW-removal-Serre}), from
Skinner--Wiles nearly-ordinary deformation theory, and from
Arthur--Clozel solvable base change used directly for the dihedral
sector. Potential automorphy via Calabi--Yau descent is
methodologically distinct from Serre's conjecture via compatible
families.

\paragraph{Sign and convention consistency.} Automorphic weight:
parallel weight \(2\) on \(F\) corresponds to weight \(2\) on
\(\mathbb{Q}\) after descent (Gee~2006 Compositio~Math.~142 \S 2).
\(L\)-parameter normalisation follows Taylor~2003 \S 1.1 (local
parameters at the archimedean places of \(F\) of the prescribed
Hodge--Tate type). Completed-cohomology normalisation follows
Emerton~2011 \S 1.2. All three sources agree after the standard
weight / \(L\)-parameter correspondences.

\paragraph{Ambient category check.} The source is the icosahedral
block of
\(\mathrm{D}^{\mathrm{sph\text{-}fin}}([GL_{2}(\mathbb{Q}_{p})])_{\chi}\).
The target is \(\mathrm{IndCoh}_{\mathrm{Nilp}}\) on the icosahedral
component of the EG stack, which is a formal spectrum
\(\mathrm{Spf}\,R^{\mathrm{univ}}_{\bar\rho,\chi}\) of a complete
Noetherian local ring of dimension \(3\). Both are stable presentable
\(\infty\)-categories with a compact generator: the completed Hecke
algebra on the source side and the structure sheaf on the target side.

\paragraph{Hypothesis domain.} Odd determinant is explicit: it is the
parity condition required for Taylor~2003 Part~II Thm~A and for the
existence of complex conjugation in the icosahedral image. The
hypothesis that the solvable chain
\(\mathbb{Q}\subset F_{1}\subset\ldots\subset F\) consists of cyclic
layers of prime degree not equal to \(p\) is explicit in
\Cref{v4-arith:w10-f:lem:BLGG-descent}. The residual failure at
non-Arthur--Clozel-cyclic layers is
\Cref{v4-arith:w10-f:conj:ns-defect}. (TW1) \(p\geq 5\) is retained.

\paragraph{Hecke-action compatibility.} The Hecke-eigenvalue match on
the icosahedral sector is conditional on the chain of \(R=T\) and
base-change identifications. Taylor gives the Artin potential
automorphy input, BLGG gives descent under its automorphic-lifting
hypotheses, and Emerton gives completed-cohomology patching under its
local--global hypotheses. At each conditional step, the Hecke action on
the automorphic side corresponds to the Frobenius trace on the Galois
side.

\paragraph{Inputs.} The automorphy and patching inputs on
\(S^{A_{5},\mathrm{odd}}\) are conditional as stated in
\Cref{v4-arith:w10-f:lem:BLGG-descent} and
\Cref{v4-arith:w10-f:lem:Emerton-patching}.
The categorical theorem also requires the DEG--completed-cohomology
generator comparison. The local--global
defect residual is \Cref{v4-arith:w10-f:conj:ns-defect}, of domain
non-formal theorem.

\paragraph{Numerical comparison.} No numerical comparison backs F1.a.1.

The six checks isolate the odd-determinant icosahedral reduction under
its stated automorphy, patching, and generator-comparison inputs. They
do not remove these inputs.

\section{Combined F1.a.1 Conditional Reduction on Five Strata}
\label{v4-arith:w10-f:combined}

\begin{theorem}[F1.a.1 equivalence on five strata, conditional on generator comparisons]
\label{v4-arith:w10-f:thm:f1a1-maximal-w10f}
\ClaimStatusNeedsVerification. For every prime \(p\geq 5\) and every central
character \(\chi\in\mathrm{Hom}(\mathbb{Z}_{p}^{\times},k^{\times})\),
assume the generator-comparison inputs in
\Cref{v4-arith:w10-f:thm:f1a1-ss-FL} and
\Cref{v4-arith:w10-f:thm:A5-odd-closure}. Then
the DEG functor \(\Phi^{\mathrm{DEG}}_{p,\chi}\) is an equivalence on
the union of:
\begin{enumerate}
\item the unramified ordinary summand (Chapter~\ref{v4-arith:w5-e:chapter});
\item the potentially Barsotti--Tate non-ordinary summand
(\Cref{v4-arith:w6-e:thm:f1a1-pBT-nonord});
\item the dihedral, nearly-ordinary, and Serre-modular non-TW sectors
(\Cref{v4-arith:w6-e:thm:TW-removal-union});
\item the Fontaine--Laffaille supersingular stratum
(\Cref{v4-arith:w10-f:thm:f1a1-ss-FL});
\item the odd-determinant icosahedral sector
(\Cref{v4-arith:w10-f:thm:A5-odd-closure}).
\end{enumerate}
This is a conditional direct-sum statement. The remaining obstructions are:
the two generator comparisons above,
the small-weight non-Fontaine--Laffaille supersingular tail
(\Cref{v4-arith:w10-f:conj:ss-small-weight}) and the
local--global defect icosahedral sub-sector
(\Cref{v4-arith:w10-f:conj:ns-defect}).
\end{theorem}

\begin{proof}
Take the direct sum of the five stated theorems. Orthogonality of the corresponding
residual-index blocks on both sides
(\Cref{v4-arith:deg-f1a:prop:spectral-blocks}) ensures the direct-sum
functor is an equivalence on the union.
\end{proof}

\begin{proposition}[F1.a.1 obstruction decomposition after Chapter~\ref{v4-arith:w10-f:chapter}]
\label{v4-arith:w10-f:prop:F1a1-table-w10f}
\ClaimStatusProvedHere. After the reductions of this chapter, the
F1.a.1 residual has four visible sub-cores:
\begin{enumerate}
\item \textbf{F1.a.1.ss-generator-match}: the DEG--Breuil--M{\'e}zard
generator comparison on the Fontaine--Laffaille supersingular stratum.
\item \textbf{F1.a.1.ss-small-weight}
(\Cref{v4-arith:w10-f:conj:ss-small-weight}): essential surjectivity
at Hodge--Tate weights \(k\in\{p-1,p\}\) supersingular boundary.
\item \textbf{F1.a.1.ns-generator-match}: the DEG--completed-cohomology
generator comparison on the odd-determinant icosahedral sector.
\item \textbf{F1.a.1.ns-defect}
(\Cref{v4-arith:w10-f:conj:ns-defect}): essential surjectivity on
the local--global descent defect icosahedral sub-sector.
\end{enumerate}
No numerical compression is recorded unless the two generator
comparisons are proved.
\end{proposition}

\begin{proof}
It follows from \Cref{v4-arith:w10-f:cor:ss-core-reduction} and
\Cref{v4-arith:w10-f:cor:ns-core-reduction}. Each corollary is
conditional on a generator comparison, hence the generator comparison
itself remains in the decomposition.
\end{proof}

\begin{proposition}[Full F1 obstruction decomposition after Chapter~\ref{v4-arith:w10-f:chapter}]
\label{v4-arith:w10-f:prop:F1-table-w10f}
\ClaimStatusProvedHere. Combining
\Cref{v4-arith:w10-f:prop:F1a1-table-w10f} with the unchanged
F1.a.2 and Eisenstein residuals, and with the corrected
\(\mathrm{L}_{\mathrm{ACD}}\) decomposition of
\Cref{v4-arith:lacd:cor:LACD-residual}, the full F1 obstruction decomposition
contains:
\begin{enumerate}
\item the four F1.a.1 sub-cores of
\Cref{v4-arith:w10-f:prop:F1a1-table-w10f};
\item \textbf{F1.a.2.class-match} (non-formal theorem; unchanged from
\Cref{v4-arith:w6-e:prop:F1-table-combined}).
\item \textbf{\(\mathrm{L}_{\mathrm{ACD}}\).Eisenstein} (non-formal theorem).
\item the local categorical comparison, the \(\infty\)-categorical
Bernstein lift, the archimedean boundary comparison, the global
identification, and (BC-diamond-glue).
\end{enumerate}
The old five-row domain estimate is therefore not a closure table.
\end{proposition}

\begin{proof}
Tabulation from \Cref{v4-arith:w10-f:prop:F1a1-table-w10f}
and \Cref{v4-arith:lacd:cor:LACD-residual}. The latter replaces the
old single BC-diamond-glue row by its comparison-data table.
\end{proof}

The two F1.a.1 residuals isolated in
Chapter~\ref{v4-arith:w6-e:chapter} reduce on their dominant sub-loci:
the Fontaine--Laffaille supersingular stratum
(\Cref{v4-arith:w10-f:thm:f1a1-ss-FL}) and the odd-determinant
icosahedral sub-sector (\Cref{v4-arith:w10-f:thm:A5-odd-closure}).
The remaining inputs are the two generator comparisons, the
small-weight Fontaine--Laffaille boundary
(\Cref{v4-arith:w10-f:conj:ss-small-weight}), and the local--global
cyclic descent defect (\Cref{v4-arith:w10-f:conj:ns-defect}). The six
assembled sources are independent from those used in
Chapter~\ref{v4-arith:w6-e:chapter}.
